\Article{7}
Putting the value of the integral \(\int xx\varphi x.\mathrm{d}x\) extended from \(x=-\infty\) to \(x=+\infty\) \(=mm\), we shall call the quantity \(m\) the \textit{mean error to be feared}, or simply the \textit{mean error}, of observations whose indeterminate errors \(x\) have the relative probability \(\varphi x\). We shall not restrict this name to direct observations, but shall also extend it to any determinations derived from observations. Care must be taken, however, not to confuse the mean error with the arithmetic mean of all the errors, of which we spoke in art. 5.

\SourceLeaf{013}Suppose several kinds of observations, or several determinations obtained from observations, to which the same precision cannot be assigned, are being compared. Their relative \textit{weight} for us will be a quantity inversely proportional to \(mm\), while their \textit{precision} is taken as simply inversely proportional to \(m\). In order, therefore, that a weight may be expressed by a number, the weight of a particular kind of observation must have been taken as the unit.

\Article{8}
If the errors of observations contain a constant part, removing it diminishes the mean error and increases the weight and precision. Retaining the notation of art. 5, and denoting the mean error of the corrected observations by \(m'\), we shall have
\[
\begin{aligned}
m'm'&=\int x'x'\varphi'x'.\mathrm{d}x'=\int(x-k)^2\varphi x.\mathrm{d}x\\
&=\int xx\varphi x.\mathrm{d}x-2k\int x\varphi x.\mathrm{d}x+kk\int\varphi x.\mathrm{d}x\\
&=mm-2kk+kk=mm-kk.
\end{aligned}
\]
If, however, instead of the true constant part \(k\), another quantity \(l\) had been subtracted from the observations, the square of the new mean error would become \(=mm-2kl+ll=m'm'+(l-k)^2\).

\Article{9}
If \(\lambda\) denotes a definite coefficient, and \(\mu\) the value of the integral \(\int\varphi x.\mathrm{d}x\) from \(x=-\lambda m\) to \(x=+\lambda m\), then \(\mu\) will be the probability that the error of an observation is less than \(\lambda m\) (without regard to sign), and \(1-\mu\) the probability of an error greater than \(\lambda m\). If, therefore, the value \(\mu=\tfrac12\) corresponds to the value \(\lambda m=\rho\), the error can fall below \(\rho\) just as readily as above \(\rho\), so that \(\rho\) can conveniently be called the \textit{probable} error. The relation between the quantities \(\lambda\), \(\mu\) evidently depends on the nature of the function \(\varphi x\), which is usually unknown. It will therefore be worthwhile to consider this relation more closely in certain special cases.

I. If the limits of all possible errors are \(-a\) and \(+a\), and all errors within these limits are equally probable, \(\varphi x\) will be constant between the limits \(x=-a\) and \(x=+a\), and consequently \(=\tfrac{1}{2a}\). Hence \(m=a\BareRoot{\tfrac13}\), and also \(\mu=\lambda\BareRoot{\tfrac13}\), as long as \(\lambda\) is not greater than \(\BareRoot{3}\); finally, \(\rho=m\BareRoot{\tfrac34}=0{,}8660254m\), and the probability that an error no greater than the mean error occurs will be \(=\BareRoot{\tfrac13}=0{,}5773503\).

II. If, as before, \(-a\) and \(+a\) are the limits of the possible errors, and the probability of the errors themselves is supposed to decrease on both sides from the error \(0\) in arithmetic progression, we shall have
\SourceLeaf{014}
\[
\begin{aligned}
\varphi x&=\frac{a-x}{aa},\quad\text{for values of }x\text{ between }0\text{ and }+a\\[3pt]
\varphi x&=\frac{a+x}{aa},\quad\text{for values of }x\text{ between }0\text{ and }-a
\end{aligned}
\]
From this we deduce \(m=a\BareRoot{\tfrac16}\), \(\mu=\lambda\BareRoot{\tfrac23}-\tfrac16\lambda\lambda\), as long as \(\lambda\) is between \(0\) and \(\BareRoot{6}\); finally, \(\lambda=\BareRoot{6}-\BareRoot{(6-6\mu)}\), as long as \(\mu\) is between \(0\) and \(1\), and consequently
\[
\rho=m(\BareRoot{6}-\BareRoot{3})=0{,}7174389m
\]
The probability of an error not exceeding the mean error will in this case be
\[
=\BareRoot{\tfrac32}-\tfrac16=0{,}6498299
\]
III. If we put the function \(\varphi x\) proportional to \(e^{-\frac{xx}{hh}}\) (which can indeed be only approximately true in nature), we must have
\[
\varphi x=\frac{e^{-\frac{xx}{hh}}}{h\BareRoot{\pi}}
\]
where \(\pi\) denotes the semicircumference of a circle of radius \(1\), whence we further deduce
\[
m=h\BareRoot{\tfrac12}
\]
(See \textit{Disquis. generales circa seriem infinitam} etc. art. 28). Further, if the value of the integral
\[
\frac{2}{\BareRoot{\pi}}\int e^{-zz}\mathrm{d}z
\]
beginning at \(z=0\) is denoted by \(\Theta z\), we shall have
\[
\mu=\Theta(\lambda\BareRoot{\tfrac12})
\]
\Needspace{14\baselineskip}
The following table gives some values of this quantity:
\begin{center}
\begin{tabular}{r|l}
\multicolumn{1}{c|}{\(\lambda\)}&\multicolumn{1}{c}{\(\mu\)}\\\hline
\(0{,}6744897\)&\(0{,}5\)\\
\(0{,}8416213\)&\(0{,}6\)\\
\(1{,}0000000\)&\(0{,}68268\mathord{\cdot}95\)\\
\(1{,}0364334\)&\(0{,}7\)\\
\(1{,}2815517\)&\(0{,}8\)\\
\(1{,}6448537\)&\(0{,}9\)\\
\(2{,}5758293\)&\(0{,}99\)\\
\(3{,}2918301\)&\(0{,}999\)\\
\(3{,}8905940\)&\(0{,}9999\)\\
\(\infty\)&\(1\)
\end{tabular}
\end{center}

\clearpage
